\section{Generating functions, branching and renewal}\label{app:C.11}

A \SI{100}{eV} electron in silicon ends as about 27 electron--hole pairs, each made by a carrier that was itself made by an earlier one. \Cref{sec:18.1,sec:18.2} treat this cascade as a branching process and read the mean energy per pair $\varepsilon_{\mathrm{eh}}$ and the Fano factor $F_{\mathrm{eh}}$ off its first two moments. Probability generating functions turn sums of independent counts into products. The branching recursion builds the statistics of a whole tree from the statistics of one step. Renewal theory says why the mean pair number becomes linear in the deposited energy, and how much of the energy is left over at the end.

\subsection{Probability generating functions}

For a random count $N = 0, 1, 2, \dots$ with probabilities $p_n$, the \emph{probability generating function} is
\begin{equation}
  \mathfrak{G}_N(s) = \langle s^N\rangle = \sum_{n\ge0} p_n\,s^n, \qquad 0 \le s \le 1 .
  \label{eq:C.11-pgf}
\end{equation}
It is a power series with nonnegative coefficients, so it is increasing and convex on $[0, 1]$, with $\mathfrak{G}_N(0) = p_0$ and $\mathfrak{G}_N(1) = 1$. Derivatives at $s = 1$ give factorial moments: $\mathfrak{G}_N'(1) = \langle N\rangle$ and $\mathfrak{G}_N''(1) = \langle N(N - 1)\rangle$. Hence
\begin{equation}
  \operatorname{Var}N = \mathfrak{G}_N''(1) + \mathfrak{G}_N'(1) - \mathfrak{G}_N'(1)^2 .
  \label{eq:C.11-var}
\end{equation}
The Poisson distribution of \cref{eq:1.5-poisson} with mean $\bar N$ has $\mathfrak{G}(s) = \sum_n\ee^{-\bar N}(\bar Ns)^n/n! = \ee^{\bar N(s - 1)}$. Then $\mathfrak{G}' (1)= \bar N$, $\mathfrak{G}''(1) = \bar N^2$, and $\operatorname{Var}N = \bar N$: the Fano factor $\operatorname{Var}N/\langle N\rangle$ of a Poisson count is 1, the reference against which the value 0.1--0.15 of silicon (\cref{sec:1.5}) is sub-Poissonian.

If $N_1$ and $N_2$ are independent, $\langle s^{N_1 + N_2}\rangle = \langle s^{N_1}\rangle\langle s^{N_2}\rangle$: the generating function of a sum is the product. If $N$ items are each kept independently with probability $f$, the kept number $N_{\mathrm{k}}$ is a sum of $N$ independent Bernoulli counts, each with generating function $1 - f + fs$. At fixed $N$, $N_{\mathrm{k}}$ has the generating function $(1 - f + fs)^N$, and averaging over $N$,
\begin{equation}
  \mathfrak{G}_{N_{\mathrm{k}}}(s) = \mathfrak{G}_N(1 - f + fs) .
  \label{eq:C.11-thin}
\end{equation}
By the chain rule $\langle N_{\mathrm{k}}\rangle = f\langle N\rangle$ and $\mathfrak{G}_{N_{\mathrm{k}}}''(1) = f^2\mathfrak{G}_N''(1) = f^2(\operatorname{Var}N - \langle N\rangle + \langle N\rangle^2)$. Inserting $\langle N_{\mathrm{k}}\rangle$ and $\mathfrak{G}_{N_{\mathrm{k}}}''(1)$ into \cref{eq:C.11-var} and dividing by $f\langle N\rangle$ gives the Fano factor of the kept items,
\begin{equation}
  F_{\mathrm{k}} = f\,F_N + 1 - f, \qquad F_N = \operatorname{Var}N/\langle N\rangle .
  \label{eq:C.11-thinfano}
\end{equation}
Random loss moves any Fano factor towards the Poisson value 1. By \cref{eq:C.11-thinfano}, a detector that collects a fraction $f_{\mathrm{det}}$ of the pairs measures $f_{\mathrm{det}}F_{\mathrm{eh}} + 1 - f_{\mathrm{det}}$.

\subsection{The branching recursion and extinction}

In the simplest branching process, the Galton--Watson process, each individual of generation $n$ has a random number of offspring, independently of all others and with one distribution, of generating function $\mathfrak{g}(s)$ and mean $\bar m = \mathfrak{g}'(1)$. Let $N_n$ be the size of generation $n$, with $N_0 = 1$, and $\mathfrak{G}_n$ its generating function. Given $N_n$, the next generation is a sum of $N_n$ independent offspring counts, so $\langle s^{N_{n+1}}\,|\,N_n\rangle = \mathfrak{g}(s)^{N_n}$, and averaging over $N_n$,
\begin{equation}
  \mathfrak{G}_{n+1}(s) = \mathfrak{G}_n\bigl(\mathfrak{g}(s)\bigr) = \mathfrak{g}\bigl(\mathfrak{G}_n(s)\bigr), \qquad \mathfrak{G}_0(s) = s .
  \label{eq:C.11-gw}
\end{equation}
The first form is what we derived; it makes $\mathfrak{G}_n$ the $n$-fold composition of $\mathfrak{g}$ with itself, so the second follows. The second form is a \emph{first-step} statement: condition on the offspring of the ancestor, and each child starts an independent copy of the process with $n$ generations to go. Differentiating at $s = 1$ gives $\langle N_n\rangle = \bar m^{\,n}$.

The probability that the line has died out by generation $n$ is $\mathfrak{G}_n(0)$. By \cref{eq:C.11-gw}, $\mathfrak{G}_{n+1}(0) = \mathfrak{g}(\mathfrak{G}_n(0))$. Since $\mathfrak{g}$ is increasing and $\mathfrak{G}_1(0) = \mathfrak{g}(0) \ge 0 = \mathfrak{G}_0(0)$, the sequence increases, and being bounded by 1 it converges to a limit $P_{\mathrm{ext}}$ with $\mathfrak{g}(P_{\mathrm{ext}}) = P_{\mathrm{ext}}$, by continuity of $\mathfrak{g}$. The equation $\mathfrak{g}(s) = s$ can have two roots in $[0, 1]$, and the iteration selects the smaller: for any root $s_*$, $0 \le s_*$ implies $\mathfrak{G}_n(0) \le \mathfrak{G}_n(s_*) = s_*$ for all $n$, so $P_{\mathrm{ext}} \le s_*$. Hence
\begin{keyresult}[Extinction probability]
The extinction probability $P_{\mathrm{ext}}$ is the smallest root in $[0, 1]$ of $\mathfrak{g}(s) = s$. It is 1 for $\bar m \le 1$ (excluding the trivial $\mathfrak{g}(s) = s$) and smaller than 1 for $\bar m > 1$~\cite[Chap.~I]{harris1963theory}.
\end{keyresult}
The dependence on $\bar m$ follows from convexity. The function $\mathfrak{g}(s) - s$ is convex, vanishes at $s = 1$ and has the slope $\bar m - 1$ there. For $\bar m \le 1$ a convex function with nonpositive slope at its zero $s = 1$ lies above its tangent there, $(\bar m - 1)(s - 1) \ge 0$ on $[0, 1)$. A root $s_0 < 1$ would therefore need $\bar m = 1$. Convexity would then hold $\mathfrak{g}(s) - s$ between the zero tangent and the zero chord from $s_0$ to 1, so it would vanish on $[s_0, 1]$. A power series that vanishes on an interval vanishes identically, and $\mathfrak{g}(s) = s$ is the trivial case excluded above. For $\bar m > 1$ the function is negative just below $s = 1$ and equals $p_0 \ge 0$ at $s = 0$, so a root lies in $[0, 1)$. For a concrete case let each individual have no offspring with probability $1 - p$ and two with probability $p$, $\mathfrak{g}(s) = 1 - p + ps^2$. The roots of $ps^2 - s + 1 - p = 0$ are $s = 1$ and $s = (1 - p)/p$. For $p \le \tfrac12$ the smaller root is 1 and the line dies out with certainty; for $p = 0.6$ it survives with probability $1 - \tfrac23 = \tfrac13$, and iterating \cref{eq:C.11-gw} from $s = 0$ converges to $\tfrac23$, the smaller root.

In a cascade, the fixed-point equation of a branching process has more solutions than the physical one, and the physical solution is the \emph{minimal} one, obtained by iterating from the terminal data. In a carrier cascade every collision lowers the energy of the particles it produces, every branch stops below the impact-ionization threshold, and termination has probability 1; the recursion is then solved from low energy upward and the solutions other than the physical one never arise.

\subsection{Branching with types and deposits}

A cascade particle carries more than a generation index: a type (electron or hole) and an energy $x$. Let a particle in state $x$ undergo its next collision in channel $c$ with probability $p_c(x)$. The collision produces particles in states $x'_1, \dots, x'_k$, drawn from a distribution that depends on $x$ and $c$, and leaves an additive deposit $d_c$: a pair, an emitted phonon energy, or a residual kinetic energy at the end. Let $Y$ be the total deposit of the tree started by the particle, and $\mathfrak{G}_x(s) = \langle s^Y\,|\,x\rangle$ for an integer deposit. If the subtrees of the $k$ products evolve independently, the first-step argument that gave \cref{eq:C.11-gw} gives
\begin{equation}
  \mathfrak{G}_x(s) = \sum_c p_c(x)\,\Bigl\langle s^{d_c}\prod_{j=1}^{k}\mathfrak{G}_{x'_j}(s)\Bigr\rangle_c ,
  \label{eq:C.11-functional}
\end{equation}
with $\langle\cdots\rangle_c$ the average over the products of channel $c$. Moments follow directly from $Y = d_c + Y_1 + \dots + Y_k$, without differentiating. Write $\bar Y(x) = \langle Y\,|\,x\rangle$ and $\overline{Y^2}(x) = \langle Y^2\,|\,x\rangle$. Averaging $Y$ and $Y^2$ over the first collision, and using $\langle Y_iY_j\rangle = \bar Y(x'_i)\bar Y(x'_j)$ for $i \ne j$ (independent subtrees),
\begin{align}
  \bar Y(x) &= \sum_c p_c(x)\Bigl\langle d_c + \sum_j\bar Y(x'_j)\Bigr\rangle_c ,
  \label{eq:C.11-mean}\\
  \overline{Y^2}(x) &= \sum_c p_c(x)\Bigl\langle d_c^2 + 2d_c\sum_j\bar Y(x'_j) + \sum_j\overline{Y^2}(x'_j) + 2\sum_{i<j}\bar Y(x'_i)\bar Y(x'_j)\Bigr\rangle_c .
  \label{eq:C.11-second}
\end{align}
\Cref{eq:C.11-mean,eq:C.11-second} hold for any additive deposit, integer or not. Every term on the right refers to products of lower energy than $x$, so a single sweep upward in $x$ solves both equations, starting from the terminal states, where $\bar Y$ and $\overline{Y^2}$ are known. The variance is $\overline{Y^2} - \bar Y^2$, and covariances of two deposits follow from the second moment of their sum.

\subsection{Renewal theory}

For deposits large compared with the gap, the mean pair number is linear in the deposited energy. Strip the cascade to a single chain: a quantity $E$, the energy, is reduced by independent losses $X_1, X_2, \dots$ with one distribution $F_{\mathrm{loss}}$ on $x > 0$, mean $\langle X\rangle$ and second moment $\langle X^2\rangle$. Let $N(E)$ be the number of losses that fit, the largest $n$ with $S_n = X_1 + \dots + X_n \le E$, and $\bar N(E) = \langle N(E)\rangle$. Condition on the first loss $X_1 = x$: either $x > E$ and $N = 0$, or the chain restarts with $E - x$ and one more count. The two cases give the \emph{renewal equation}
\begin{equation}
  \bar N(E) = F_{\mathrm{loss}}(E) + \int_0^E\bar N(E - x)\,\dd F_{\mathrm{loss}}(x) .
  \label{eq:C.11-renewal}
\end{equation}
Let $Z(E) = z(E) + \int_0^E Z(E - x)\,\dd F_{\mathrm{loss}}(x)$, with $z$ nonnegative, nonincreasing and integrable (more generally, directly Riemann integrable).
\begin{keyresult}[Key renewal theorem]
If $F_{\mathrm{loss}}$ is not concentrated on a lattice $\{0, d, 2d, \dots\}$ and has finite mean, then
\begin{equation}
  \lim_{E\to\infty} Z(E) = \frac{1}{\langle X\rangle}\int_0^\infty z(y)\,\dd y .
  \label{eq:C.11-key}
\end{equation}
The proof is in Ref.~\cite[Chap.~XI.1]{feller1971introduction}.
\end{keyresult}

We first apply the theorem to the energy left over at the end, the \emph{undershoot} $E_{\mathrm{res}}(E) = E - S_{N(E)}$, the part of $E$ that no further loss fits into. Fix $y \ge 0$ and let $Z(E) = P\bigl(E_{\mathrm{res}}(E) > y\bigr)$. If the first loss already exceeds $E$, the residual is $E$ itself, which exceeds $y$ when $E > y$; otherwise the chain restarts with $E - x$. So $Z = z + Z * F_{\mathrm{loss}}$ with $z(E) = \Theta(E - y)\,[1 - F_{\mathrm{loss}}(E)]$, nonincreasing beyond $y$. By \cref{eq:C.11-key},
\begin{equation}
  \lim_{E\to\infty}P\bigl(E_{\mathrm{res}} > y\bigr) = \frac{1}{\langle X\rangle}\int_y^\infty\bigl[1 - F_{\mathrm{loss}}(x)\bigr]\dd x ,
  \qquad
  \langle E_{\mathrm{res}}\rangle \to \frac{\langle X^2\rangle}{2\langle X\rangle} .
  \label{eq:C.11-undershoot}
\end{equation}
The limiting density of the residual is $[1 - F_{\mathrm{loss}}(y)]/\langle X\rangle$, and the mean residual follows by integrating $y$ against this density by parts: $\int_0^\infty y[1 - F_{\mathrm{loss}}(y)]\,\dd y = \tfrac12\int_0^\infty y^2\,\dd F_{\mathrm{loss}}(y) = \tfrac12\langle X^2\rangle$. A residual is not a typical loss: long losses cover more of the energy axis, so the gap in which the chain ends is biased towards them.

The mean count follows from energy bookkeeping. With the convention $S_0 = 0$, the loss that first overshoots $E$ is $X_{N+1}$, and $N + 1$ is a stopping time: whether it has occurred is decided by $X_1, \dots, X_{N+1}$. Wald's identity then gives $\langle S_{N+1}\rangle = \langle X\rangle(\bar N + 1)$~\cite{feller1971introduction}. The overshoot $S_{N+1} - E$ obeys a renewal equation of the same kind: if the first loss exceeds $E$, the overshoot is $X_1 - E$, which exceeds $y$ with probability $1 - F_{\mathrm{loss}}(E + y)$; otherwise the chain restarts. With $z(E) = 1 - F_{\mathrm{loss}}(E + y)$, \cref{eq:C.11-key} gives the limit $\langle X\rangle^{-1}\int_y^\infty[1 - F_{\mathrm{loss}}(x)]\,\dd x$, the law of \cref{eq:C.11-undershoot}, with mean $\langle X^2\rangle/(2\langle X\rangle)$. Hence
\begin{equation}
  \bar N(E) = \frac{E}{\langle X\rangle} + \frac{\langle X^2\rangle}{2\langle X\rangle^2} - 1 + o(1) .
  \label{eq:C.11-linear}
\end{equation}
The leading term is the linear yield, $\bar N(E)/E \to 1/\langle X\rangle$; in the cascade the mean loss per pair is $\varepsilon_{\mathrm{eh}}$, and \cref{sec:18.2} derives $\bar n(\omega) = \omega/\varepsilon_{\mathrm{eh}} + o(\omega)$ in this way. The fluctuations follow from the central limit theorem for $S_n$. $N(E) < n$ means $S_n > E$, and $S_n$ is approximately Gaussian with mean $n\langle X\rangle$ and variance $n\sigma_X^2$, $\sigma_X^2 = \langle X^2\rangle - \langle X\rangle^2$. Setting $n = E/\langle X\rangle + u\,(E\sigma_X^2/\langle X\rangle^3)^{1/2}$ and keeping the leading order in $E$ gives $P(N < n) = P(S_n > E) \to \Phi(u)$, the standard normal distribution function. So $N(E)$ is Gaussian at large $E$ with
\begin{equation}
  \operatorname{Var}N(E) \simeq \frac{\sigma_X^2}{\langle X\rangle^3}\,E, \qquad
  \frac{\operatorname{Var}N}{\bar N} \to \frac{\sigma_X^2}{\langle X\rangle^2} .
  \label{eq:C.11-fano}
\end{equation}
The Fano factor of a single renewal chain is the squared relative width of one loss. A deterministic loss gives zero, and an exponential loss, for which $\sigma_X = \langle X\rangle$, gives the Poisson value 1. For losses uniform on $[0, 1]$, $\langle X\rangle = \tfrac12$ and $\langle X^2\rangle = \tfrac13$. \Cref{eq:C.11-undershoot,eq:C.11-linear,eq:C.11-fano} then give $\langle E_{\mathrm{res}}\rangle = \tfrac13$, $\bar N - 2E = -\tfrac13$ and a Fano factor $\tfrac13$. A numerical solution of \cref{eq:C.11-renewal} at $E = 20$ reproduces the offset to $10^{-6}$, and a simulation of $2\times10^{5}$ chains at $E = 50$ reproduces $\langle E_{\mathrm{res}}\rangle$ and the Fano factor to \SI{1}{\percent}.

Each hypothesis of the key renewal theorem fails somewhere in a cascade. A stationary loss law, independent of $E$, fails when the losses depend on the carrier energy; the theorem then applies only to the high-energy part of the chain. A non-lattice law fails when only a few loss values occur. For a deterministic loss $d$, $N(E) = \lfloor E/d\rfloor$ and $\bar N - E/d = -\{E/d\}$, the negative fractional part, which oscillates forever between $-1$ and 0. Its average over a period, $-\tfrac12$, is what \cref{eq:C.11-linear} predicts with $\langle X^2\rangle = d^2$, but the limit does not exist. A model whose losses are integer combinations of a gap and one phonon energy has a lattice loss law, and its yield carries a periodic term. A finite second moment fails for a fast primary that loses energy in knock-on collisions with a tail $\dd F_{\mathrm{loss}}/\dd x \propto x^{-2}$ up to $E/2$. Then $\langle X^2\rangle$ grows in proportion to $E$, the residual of \cref{eq:C.11-undershoot} grows without bound, and the corrections to linearity can be logarithmic in $E$ rather than constant.
